\documentclass[10pt, aspectratio=169]{beamer}
% Preámbulo modular para presentaciones Beamer (Modelación Estadística)
% Diseño y paleta institucional del Tecnológico de Monterrey

\documentclass[aspectratio=169,xcolor={dvipsnames,table}]{beamer}

\usepackage[utf8]{inputenc}
\usepackage[spanish,mexico]{babel}
\usepackage{amsmath,amssymb,amsfonts,mathtools}
\usepackage{graphicx}
\usepackage{listings}
\usepackage{booktabs}
\usepackage{tikz}
\usetikzlibrary{matrix,arrows,decorations.pathmorphing,shapes.geometric,positioning}

% Paleta Institucional Tecnológico de Monterrey
\definecolor{TecAzul}{HTML}{0039A6}       % Azul TEC: color primario institucional
\definecolor{TecAzulOscuro}{HTML}{1A2E51} % Azul marino oscuro
\definecolor{TecRojo}{HTML}{EC2661}       % Rojo de acento (paleta miTec)
\definecolor{TecGris}{HTML}{646464}       % Gris neutro para comentarios y subtítulos
\definecolor{TecFondoBloque}{HTML}{F4F6F9}% Fondo suave para bloques

% Configuración de Tema Beamer
\usetheme{Boadilla}
\usecolortheme[named=TecAzulOscuro]{structure}
\setbeamercolor{alerted text}{fg=TecRojo}
\setbeamercolor{palette primary}{bg=TecAzulOscuro,fg=white}
\setbeamercolor{palette secondary}{bg=TecAzul,fg=white}
\setbeamercolor{palette tertiary}{bg=TecAzulOscuro!80!black,fg=white}
\setbeamercolor{title}{fg=TecAzulOscuro,bg=TecFondoBloque}
\setbeamercolor{frametitle}{fg=TecAzulOscuro,bg=TecFondoBloque}

% Bloques personalizados elegantes
\setbeamercolor{block title}{bg=TecAzulOscuro,fg=white}
\setbeamercolor{block body}{bg=TecFondoBloque,fg=black}
\setbeamercolor{block title alerted}{bg=TecRojo,fg=white}
\setbeamercolor{block body alerted}{bg=TecRojo!8,fg=black}
\setbeamercolor{block title example}{bg=TecAzul,fg=white}
\setbeamercolor{block body example}{bg=TecAzul!8,fg=black}

% Redondear bordes y añadir sombra sutil
\setbeamertemplate{blocks}[rounded][shadow=true]
\setbeamertemplate{navigation symbols}{}

% Pie de página limpio con número de diapositiva
\setbeamertemplate{footline}{
  \leavevmode%
  \hbox{%
  \begin{beamercolorbox}[wd=.7\paperwidth,ht=2.25ex,dp=1ex,left]{author in head/foot}%
    \hspace*{2ex}\insertshortauthor~~(\insertshortinstitute) \hfill \insertshorttitle
  \end{beamercolorbox}%
  \begin{beamercolorbox}[wd=.3\paperwidth,ht=2.25ex,dp=1ex,right]{date in head/foot}%
    \insertshortdate{}\hspace*{2em}
    \insertframenumber{} / \inserttotalframenumber\hspace*{2ex}
  \end{beamercolorbox}}%
  \vskip0pt%
}

% Configuración de Listings de Python para visualización pedagógica de código
\lstset{
    language=Python,
    basicstyle=\ttfamily\footnotesize,
    keywordstyle=\color{TecAzulOscuro}\bfseries,
    stringstyle=\color{TecRojo},
    commentstyle=\color{TecGris}\itshape,
    showstringspaces=false,
    breaklines=true,
    frame=lines,
    rulecolor=\color{TecAzulOscuro!30},
    backgroundcolor=\color{TecFondoBloque},
    aboveskip=0.5em,
    belowskip=0.5em,
    literate={á}{{\'a}}1 {é}{{\'e}}1 {í}{{\'i}}1 {ó}{{\'o}}1 {ú}{{\'u}}1
             {Á}{{\'A}}1 {É}{{\'E}}1 {Í}{{\'I}}1 {Ó}{{\'O}}1 {Ú}{{\'U}}1
             {ñ}{{\~n}}1 {Ñ}{{\~N}}1 {¿}{{?`}}1 {¡}{{!`}}1
}

% Metadatos institucionales por defecto
\title[Modelación Estadística]{Modelación Estadística para Ciencia de Datos e Ingeniería}
\author[J. Castillo Colmenares]{Juliho David Castillo Colmenares}
\institute[Tec de Monterrey]{Tecnológico de Monterrey \\ Escuela de Ingeniería y Ciencias}
\date{\today}

% Comandos y macros estadísticas compartidas para las presentaciones Beamer (Metropolis)

% Conjuntos numéricos básicos
\newcommand{\R}{\mathbb{R}}
\newcommand{\N}{\mathbb{N}}
\newcommand{\Q}{\mathbb{Q}}
\newcommand{\Z}{\mathbb{Z}}
\newcommand{\set}[1]{\left\{ #1 \right\}}
\newcommand{\abs}[1]{\left|#1\right|}
\newcommand{\norm}[1]{\left\| #1 \right\|}

% Comandos estadísticos y probabilísticos del curso
\newcommand{\Var}{\operatorname{Var}}
\newcommand{\cov}{\operatorname{Cov}}
\newcommand{\corr}{\rho}
\newcommand{\comb}[2]{\begin{pmatrix} #1 \\ #2 \end{pmatrix}}
\newcommand{\s}{\sigma}
\newcommand{\particion}{\mathfrak{P}}
\newcommand{\card}[1]{\#\left( #1 \right)}
\newcommand{\seq}[3]{\set{#1}_{#2}^{#3}}
\newcommand{\E}{\mathbb{E}}
\newcommand{\Prob}{\mathbb{P}}

% Entornos pedagógicos y cajas en español académico natural
\newenvironment{discusion}{%
  \begin{alertblock}{Pregunta de Discusión}%
}{%
  \end{alertblock}%
}

\newenvironment{reflexion}{%
  \begin{alertblock}{Reflexión para el Aula}%
}{%
  \end{alertblock}%
}

\newenvironment{paradoja}{%
  \begin{alertblock}{Paradoja de Exploración}%
}{%
  \end{alertblock}%
}

\newenvironment{motivacion}{%
  \begin{exampleblock}{Motivación: Ciencia de Datos en la Práctica}%
}{%
  \end{exampleblock}%
}

\newenvironment{retoclase}{%
  \begin{block}{Reto Rápido en Equipo}%
}{%
  \end{block}%
}


\title[Implementación con Python]{Sección 09.06: Implementación en Python con \texttt{statsmodels}}
\subtitle{De las Fórmulas Manuales a la Tabla de Resumen Profesional}
\author[J. Castillo Colmenares]{Juliho Castillo Colmenares}
\institute{Tecnológico de Monterrey}
\date{\vspace{-1.2cm}}

\begin{document}

% Slide 1: Portada
\begin{frame}[plain]
  \vspace{-0.8cm}
  \titlepage
\end{frame}

% Slide 2: Hoja de Ruta
\begin{frame}{Hoja de Ruta --- Capítulo 09: Regresiones Lineales y Múltiples}
  \scriptsize
  ¿Dónde nos encontramos en el desarrollo modular del capítulo?
  \begin{itemize}\itemsep=0.02em
    \item \textbf{09.01} Correlación como Premisa de la Regresión
    \item \textbf{09.02} Introducción a la Regresión Lineal
    \item \textbf{09.03} Matemáticas de la Regresión: Derivación MCO y $R^2$
    \item \textbf{09.04} Regresión Lineal sobre Datos Simulados
    \item \textbf{09.05} Coeficientes Óptimos, Pruebas $t$/$F$ y RSE
    \item \textbf{09.06} Implementación en Python con \texttt{statsmodels} \emph{(Hoy)}
    \item \textbf{09.07} Regresión Múltiple, Ecuación Normal y Ridge/Lasso
    \item \textbf{09.08} Validación de Modelos y $k$-fold Cross-Validation
    \item \textbf{09.09} Diagnóstico de Residuos y Multicolinealidad (VIF)
    \item \textbf{09.10} Regresión con \texttt{scikit-learn}
    \item \textbf{09.11} Variables Categóricas y Variables Muda
    \item \textbf{09.12} Transformaciones No Lineales y Regresión Polinomial
  \end{itemize}
  \pause
  \vspace{-0.05cm}
  \begin{block}{Objetivo de la Sesión}
    Reconocer cómo la librería profesional \texttt{statsmodels} condensa en una sola tabla de resumen exactamente los mismos cálculos derivados manualmente en las Secciones 09.03-09.05: coeficientes, errores estándar, estadísticos $t$/$F$, $R^2$ y RSE.
  \end{block}
\end{frame}

% Slide 3: Motivación
\begin{frame}{De las Fórmulas Manuales a la Herramienta Profesional}
  \footnotesize
  \begin{motivacion}
    Hemos derivado y calculado manualmente cada componente de un modelo de regresión. En la práctica profesional, librerías como \texttt{statsmodels} automatizan estos cálculos en una sola línea de código, produciendo una tabla de resumen estandarizada usada en la industria y la academia.
  \end{motivacion}

  \vspace{0.15cm}
  \pause
  Python ofrece dos rutas principales para regresión: \texttt{statsmodels.formula.api} (con enfoque estadístico/inferencial, similar a R) y \texttt{scikit-learn} (con enfoque de aprendizaje automático, Sección 09.10).
\end{frame}

% Slide 4: El método ols de statsmodels
\begin{frame}[fragile]{El Método \texttt{ols} de \texttt{statsmodels}}
  \footnotesize
  \begin{block}{Ajuste con Notación de Fórmula}
    \begin{verbatim}
import statsmodels.formula.api as smf
model = smf.ols(formula='Sales~TV', data=advert).fit()
print(model.params)
    \end{verbatim}
  \end{block}
  \pause

  \vspace{0.1cm}
  \begin{alertblock}{Salida}
    \texttt{Intercept} $\approx7.03$, \texttt{TV} $\approx0.0475$ --- exactamente los mismos $\hat\alpha,\hat\beta$ que calcularíamos a mano con las fórmulas de la Sección 09.03.
  \end{alertblock}
\end{frame}

% Slide 5: La tabla de resumen completa
\begin{frame}{La Tabla de Resumen \texttt{model.summary()}}
  \footnotesize
  Una sola llamada (\texttt{model.summary()}) reporta simultáneamente:
  \begin{itemize}\itemsep=0.08cm
    \item \textbf{Coeficientes} y sus \textbf{errores estándar} (Sección 09.05). \pause
    \item \textbf{Estadísticos $t$ y valores-$p$} por coeficiente (Sección 09.05). \pause
    \item \textbf{$R^2$ y $R^2$ ajustado} (Sección 09.03). \pause
    \item \textbf{Estadístico $F$} global y su valor-$p$ (Sección 09.05). \pause
    \item \textbf{Durbin-Watson}, AIC, BIC --- diagnósticos que profundizaremos en la Sección 09.09.
  \end{itemize}
\end{frame}

% Slide 6: Predicción y RSE
\begin{frame}[fragile]{Predicción y Error Estándar Residual}
  \footnotesize
  \begin{block}{Predicción sobre Nuevos Datos}
    \begin{verbatim}
sales_pred = model.predict(pd.DataFrame(advert['TV']))
    \end{verbatim}
  \end{block}
  \pause

  \vspace{0.1cm}
  \begin{alertblock}{RSE y Error Relativo}
    Con $\text{RSE}\approx3.25$ sobre una venta promedio de $\approx14.02$, el error relativo del modelo es $\approx23\%$: significativo, y motiva agregar más predictores (\texttt{Radio}, \texttt{Newspaper}) --- el tema de la Sección 09.07.
  \end{alertblock}
\end{frame}

% Slide 7: Lab Python (Bloque 1)
\begin{frame}[fragile]{Laboratorio en Python: Reconstrucción del Conjunto de Datos}
  \scriptsize
  Reconstrucción sintética del clásico ejemplo publicitario TV/Ventas (\texttt{09.06\_statsmodels\_style\_summary.py}):
  {\lstset{basicstyle=\fontsize{3.6pt}{4.3pt}\selectfont\ttfamily,aboveskip=0.05em,belowskip=0.05em}
  \lstinputlisting[firstline=18, lastline=25]{../../code/09_regresiones/09.06_statsmodels_style_summary.py}}
\end{frame}

% Slide 8: Lab Python (Bloque 2)
\begin{frame}[fragile]{Laboratorio en Python: Tabla de Resumen Completa desde Cero}
  \scriptsize
  Reproducción manual de cada entrada de una tabla \texttt{statsmodels} usando solo \texttt{numpy}/\texttt{scipy}:
  {\lstset{basicstyle=\fontsize{3.0pt}{3.7pt}\selectfont\ttfamily,aboveskip=0.05em,belowskip=0.05em}
  \lstinputlisting[firstline=29, lastline=48]{../../code/09_regresiones/09.06_statsmodels_style_summary.py}}
\end{frame}

% Slide 9: Lab Python (Bloque 3)
\begin{frame}[fragile]{Laboratorio en Python: $R^2$, Estadístico $F$ y RSE}
  \scriptsize
  Cálculo final de las métricas globales de calidad del ajuste:
  {\lstset{basicstyle=\fontsize{3.4pt}{4.1pt}\selectfont\ttfamily,aboveskip=0.05em,belowskip=0.05em}
  \lstinputlisting[firstline=51, lastline=64]{../../code/09_regresiones/09.06_statsmodels_style_summary.py}}
\end{frame}

% Slide 10: Salida Terminal
\begin{frame}[fragile]{Laboratorio en Python: Salida en Terminal}
  \scriptsize
  Salida estándar generada al ejecutar el script del laboratorio en Python:
  \vspace{0.02cm}
  \begin{block}{Salida Estándar en Consola}
    \fontsize{3.4pt}{4.1pt}\selectfont
    \begin{verbatim}
=== Block 1: Reconstructing Advertising-Style Dataset ===
n=200 observations; mean(TV)=146.74, mean(Sales)=14.05

=== Block 2: Full OLS Summary Table ===
Coefficient    Estimate     Std Err         t       P>|t|
Intercept        6.0202      0.4670     12.8914    5.24e-28
TV               0.0547      0.0028     19.7788    9.01e-49

=== Block 3: R^2, F-Statistic, and RSE ===
R-squared = 0.6640
F-statistic = 391.2003 (p-value=9.01e-49)
RSE = 3.2595, mean(Sales) = 14.0533, relative error = 0.2319 (23.2%)
    \end{verbatim}
  \end{block}
\end{frame}

% Slide 11: Comparación con el enunciado del libro
\begin{frame}{Comparación: Nuestro Motor vs. \texttt{statsmodels}}
  \footnotesize
  \begin{itemize}\itemsep=0.12cm
    \item Nuestra reconstrucción sintética produce, para esta muestra particular, intercepto $\approx6.02$, pendiente de \texttt{TV} $\approx0.055$, $R^2\approx0.664$, RSE$\approx3.26$ con error relativo $\approx23\%$ --- el mismo orden de magnitud que el ejemplo clásico del libro (intercepto $\approx7.03$, pendiente $\approx0.0475$, $R^2\approx0.61$); las diferencias específicas se deben a la variabilidad muestral inherente a una reconstrucción sintética distinta del conjunto de datos original. \pause
    \item Esto confirma que nuestro motor manual de MCO (Secciones 09.03-09.05) calcula \textbf{exactamente} lo mismo que reportaría \texttt{statsmodels} sobre los mismos datos --- la librería no usa matemática distinta, solo automatiza el cálculo. \pause
    \item El error relativo de $\approx23\%$ es sustancial: motiva incorporar más predictores en un modelo múltiple (Sección 09.07).
  \end{itemize}
\end{frame}

% Slide 12: Cierre
\begin{frame}{Cierre de Sección: Implementación con \texttt{statsmodels}}
  \begin{reflexion}
    Toda la maquinaria derivada manualmente en las Secciones 09.03-09.05 --- coeficientes MCO, pruebas $t$/$F$, RSE --- es exactamente lo que una librería profesional automatiza en una sola tabla de resumen. Entender las matemáticas subyacentes es lo que permite interpretar correctamente esa tabla, en lugar de tratarla como una caja negra.
  \end{reflexion}

  \vspace{0.2cm}
  \pause
  \begin{block}{Lecciones Clave}
    \begin{itemize}\itemsep=0.08cm
      \item \textbf{\texttt{statsmodels.formula.api.ols}:} sintaxis de fórmula estilo R, enfoque inferencial.
      \item \textbf{\texttt{model.summary()}:} condensa coeficientes, errores estándar, $t$, $F$, $R^2$, y diagnósticos en una tabla.
      \item \textbf{Ningún "truco" adicional:} la librería reproduce exactamente las fórmulas ya derivadas, no introduce matemática nueva.
    \end{itemize}
  \end{block}
\end{frame}

% Slide 13: Perspectiva Modular
\begin{frame}{Perspectiva Modular --- El Próximo Reto}
  \scriptsize
  \begin{retoclase}
    El error relativo de $\approx23\%$ de nuestro modelo con un solo predictor sugiere que \texttt{TV} no cuenta toda la historia. La Sección 09.07 generaliza el modelo a \textbf{múltiples} predictores simultáneos (\texttt{TV}, \texttt{Radio}, \texttt{Newspaper}) mediante la Ecuación Normal en notación matricial, y presenta Ridge/Lasso para cuando esos predictores están correlacionados entre sí.
  \end{retoclase}

  \vspace{0.08cm}
  \pause
  \begin{alertblock}{Preparación Recomendada}
    \begin{itemize}\itemsep=0.06cm
      \item Repasa la notación matricial básica: vectores, matrices, transpuesta e inversa.
      \item Reflexiona sobre por qué agregar predictores razonables debería, en principio, reducir el RSE.
    \end{itemize}
  \end{alertblock}
\end{frame}

\end{document}
